\Subhead{2.2.1}{Простейшие модели диспергирующих сред. Фазовая и групповая скорости}

Дисперсия — зависимость фазового набега от частоты, описываемая $k(\omega)$. Фазовая скорость $v_{\rm ph}=\omega/k$ относится к монохроматической фазе, групповая $v_g=(\dd k/\dd\omega)^{-1}$ — к огибающей узкого пакета при слабом поглощении и слабом искажении. Скорость фронта определяется причинностью и не обязана совпадать с $v_g$ в области аномальной дисперсии.

Простейшие модели: цепочка связанных осцилляторов, волновод с $k_z^2=\omega^2/c^2-k_c^2$, плазма $\varepsilon_r=1-\omega_p^2/\omega^2$ и связанный электрон Лоренца. Для последнего
\begin{equation}
 \ddot x+\Gamma\dot x+\omega_a^2x=-eE/m,
 \qquad \varepsilon_r(\omega)=1+\frac{Ne^2/(m\varepsilon_0)}{\omega_a^2-\omega^2-\ii\Gamma\omega}.
\end{equation}
Резонансная поляризация запаздывает, одновременно создавая дисперсию и поглощение.

Пакет около $\omega_0$ строится из спектра $a(\Omega)$:
\begin{equation}
 u=\operatorname{Re}\left\{\e^{\ii(k_0z-\omega_0t)}
 \int a(\Omega)\e^{-\ii\Omega t+\ii[k(\omega_0+\Omega)-k_0]z}\dd\Omega\right\}.
\end{equation}
При $k\simeq k_0+k'_0\Omega$ огибающая просто переносится как $A(t-k'_0z)$. При втором приближении появляется квадратичная спектральная фаза $k''_0z\Omega^2/2$ и разные части спектра приходят с разной групповой задержкой.
